\documentclass[sigconf]{acmart}
% Cleaning
\pagestyle{fancy} % removes running headers
\settopmatter{printacmref=false} % Removes citation information below abstract
\renewcommand\footnotetextcopyrightpermission[1]{} % removes footnote with conference information in first column
\pagestyle{plain} % removes running headers
\fancyhead{}
\settopmatter{printacmref=false, printfolios=false}
% Copyright
\setcopyright{none}
% \setcopyright{acmcopyright}
% \setcopyright{acmlicensed}
% \setcopyright{rightsretained}
% \setcopyright{usgov}
% \setcopyright{usgovmixed}
% \setcopyright{cagov}
% \setcopyright{cagovmixed}
\settopmatter{printacmref=false} % Removes citation information below abstract
\acmDOI{}
\acmConference[CSC503/SENG474]{CSC503 \texttimes{} SENG474: Data Mining}{Fall 2026}{University of Victoria}

\usepackage{booktabs} % For formal tables
\usepackage{url}
\usepackage{algorithm}
\usepackage[noend]{algpseudocode}
\usepackage{tikz}
\usetikzlibrary{arrows.meta,positioning,fit}


\setlength{\parskip}{0pt}
\setlength{\parsep}{0pt}
\setlength{\headsep}{0pt}
\setlength{\topskip}{0pt}
\setlength{\topmargin}{0pt}
\setlength{\topsep}{0pt}
\setlength{\partopsep}{0pt}
%\linespread{0.95}
\usepackage{mdwlist}
\usepackage{amsmath}

\begin{document}
\title{Evaluation of Proposed Geospatial Foundation Model Embedding-Augmented Multi-Attention-U-Net 
for Next-Day Wildfire Spread Prediction}

\author{Sebastian Brubaker}
\affiliation{\institution{University of Victoria}\city{Victoria}\state{British Columbia}\country{Canada}}
\email{sebastian.brubaker@gmail.com}

\author{Armaan Binpal}
\affiliation{\institution{University of Victoria}\city{Victoria}\state{British Columbia}\country{Canada}}
\email{armaanbinpal1@gmail.com}

\author{James Kim}
\affiliation{\institution{University of Victoria}\city{Victoria}\state{British Columbia}\country{Canada}}
\email{jknhjk03@uvic.ca}

\author{Bowen Zhang}
\affiliation{\institution{University of Victoria}\city{Victoria}\state{British Columbia}\country{Canada}}
\email{7pixels.v@gmail.com}


\begin{abstract}
This evaluation establishes the efficacy of the proposed geospatial foundation model 
embedding-augmented multi-attention U-Net for next-day wildfire spread prediction. An experimental
approach is used to benchmark the proposed method against existing methods and isolate the key factors 
that determine performance. {\color{red}[KEY FINDINGS HERE]}. 

Project repository: https://github.com/CSC503-Wildfire-Team/Wildfire-Prediction-DL
\end{abstract}

\keywords{wildfire spread prediction, geospatial foundation models, AlphaEarth embeddings, remote sensing, semantic segmentation}
\maketitle

\section{Introduction}
The proposed method comprises a custom U-Net with gated-skip-connection spatial attention \cite{almakrami} and 
Squeeze-Excitation-based channel attention \cite{SE2017} which uses the conventional environmental predictors 
of the WildfireSpreadTS dataset \cite{gerard2023wildfirespreadts} augmented with AlphaEarth Foundations (AEF) 
embeddings \cite{brown} to predict the active-fire mask at day $t+1$ given the information available at day $t$.

The efficacy of the proposed method is evaluated through an empirical approach; feature ablation isolates 
the predictive performance increase yielded by the AEF embeddings, while architecture ablation isolates the 
importance of spatial and channel attention.

\section{Research Questions}
The proposed method is evaluated based on the following research questions:

\textbf{RQ1. Does augmenting WildfireSpreadTS environmental predictors with AlphaEarth 
Foundations embeddings improve next-day active-fire prediction performance?} Two multi-attention U-Net models are 
compared; The geospatial representation-augmented model is benchmarked against a conventional environmental 
input-only model. 
The two respective models are benchmarked against the WildfireSpreadTS benchmark models 
\cite{gerard2023wildfirespreadts} as all models are developed using WildfireSpreadTS. 

\textbf{RQ2. Do channel-wise and spatial attention mechanisms improve next-day 
active-fire prediction compared with a conventional U-Net?} Four architectures are instantiated and compared: a 
conventional U-Net, a gated-skip-connection spatial attention U-Net \cite{almakrami}, a 
Squeeze-Excitation-based channel attention \cite{SE2017} U-Net, and a multi-attention U-Net. Again, additional 
comparisons are made against the WildfireSpreadTS benchmark models \cite{gerard2023wildfirespreadts}.


\section{Methodology \& Experimental Design}
{\color{red} [Explain why and how ablating isolates feature predictive power. 
Explain why baselines are used and justify each one.
EXPLAIN AND JUSTIFY THE FOLLOWING:]}
\subsection{Dataset}
\textbf{WildfireSpreadTS.} All models are trained and evaluated on WildfireSpreadTS
(WSTS)~\cite{gerard2023wildfirespreadts}. WSTS was selected for five reasons. (i) \emph{Relevance:} its
central task is exactly the one studied here, predicting the next day's binarized VIIRS active-fire map
from the observations available on the current day. (ii) \emph{Label availability:} every daily image
carries the active-fire band from which the next-day ground-truth mask is derived, so no additional
labelling is required. (iii) \emph{Comparability:} WSTS ships published benchmark models and a
year-based cross-validation protocol, so our models can be placed alongside existing results rather than
evaluated in isolation. (iv) \emph{Accessibility:} the data are openly available on Zenodo under
CC-BY-4.0 and the loader code under MIT, so the experiments can be reproduced by others.
(v) \emph{Georeferencing:} each fire event is delivered as GeoTIFFs with a coordinate reference system
and affine transform, which is a precondition for this study, since the AEF embeddings must be aligned
pixel-for-pixel with the WSTS grid.

The full dataset contains 13,607 daily images across 607 fire events larger than 1,000\,ha in the
western United States from January 2018 to October 2021, each covering a $1^\circ\times1^\circ$ area
around the fire at 375\,m resolution~\cite{gerard2023wildfirespreadts}. Each image has 23 channels:
VIIRS reflectance (M11, I2, I1), NDVI and EVI2, GRIDMET weather (precipitation, wind speed and
direction, minimum and maximum temperature, specific humidity, energy release component, PDSI), GFS
weather forecasts, topography (slope, aspect, elevation), MODIS land cover, and the active-fire
band~\cite{gerard2023wildfirespreadts}.

\textbf{AlphaEarth Foundations embeddings.} The embedding features are taken from the AlphaEarth
Foundations Satellite Embedding dataset~\cite{brown}, distributed through Google Earth Engine. It was
chosen because it is, to our knowledge, the only openly released annual, global, 10\,m geospatial
foundation-model product: every pixel is a unit-length 64-dimensional vector summarising a year of
multi-sensor observations, with Landsat, Sentinel-1 and Sentinel-2 among the inputs~\cite{brown}.

Ingesting AEF at 10\,m for all 607 events exceeded the Earth Engine compute quota available to the
project, so all experiments use a fixed subset of 40 fire events per year (160 events), taken as the
first 40 events of each year in WSTS event-ID order. The subset was fixed before any model was trained
and is shared by every model and fold, so all comparisons use identical fires and day pairs.
\subsection{Preprocessing and Data Splits}
{\color{red}[Describe and justify our preprocessing citing \cite{gerard2023wildfirespreadts}, \cite{brown} and my medoid algorithm from M3.]}
\subsection{Models and Baselines}
{\color{red}[Explain each model--\textbf{config level not architecture details}--and clearly justfy the use of our given baselines and specify 
exactly what we \textbf{can} and \textbf{cannot} infer from each ablation]}
\subsection{Training and Evaluation Metrics}
{\color{red} [Specify all training configs and randomness control--reader should be able to replicate. Explain each metric and loss 
(don't worry about justify the loss function, I did that in M3).]}

\section{Results}
{\color{red}[PRESENT/REPORT DON'T ATTRIBUTE OR DISCUSS. Add a nice looking table.]}
\subsection{RQ1: Contribution of AEF Embeddings}
\subsection{RQ2: Contribution of Attention Mechanisms}

\section{Discussion}
{\color{red}[Discuss results and draw conclusions.]}


\bibliographystyle{ACM-Reference-Format}
\bibliography{bibliography}

\end{document}